% GENERATED FILE. Edit canonical lessons, then run npm run generate:books.
% canonical-source-hash: fnv1a64:aff691ea5ad68afe
% canonical-lessons: GU-C466-kayar, GU-C466-nidar, GU-C466-saurya, GU-C466-khamir, GU-C466-dagmagvu

\chapter{Brave Hearts}
\label{ch:gu-brave-hearts}
\hlchaptermodality{466}

\begin{chapteropening}
\textbf{By the end of this chapter:} \emph{I can say who is brave and who is afraid.}

Complete the last lesson of chapter 466: I can say who is brave and who is afraid.
\end{chapteropening}

\section[\texorpdfstring{\gu{કાયર} (kāyar)}{કાયર (kāyar)}]{\gu{કાયર} (kāyar) — cowardly}

\label{lesson:GU-C466-kayar}



\begin{quote}
Before the new one: say the Gujarati for egotism, then the Gujarati for to get used to.
\end{quote}

\subsection*{You'll want to know: \gu{કાયર}}
\textbf{\gu{કાયર}} — \emph{kāyar} — \textquotedblleft{}cowardly\textquotedblright{}.

\begin{quote}
\gu{કાયર} \gu{માણસ} \gu{જોખમ} \gu{જોઈને} \gu{ભાગી} \gu{જાય} \gu{છે}.
\end{quote}

\begin{quote}
\emph{kāyar māṇas jokham joīne bhāgī jāy che.}
\end{quote}

\begin{quote}
A coward runs away at the sight of danger.
\end{quote}

\begin{cousinweb}
\textbf{\gu{કાયર}} (\emph{kāyar}) is inherited from Sanskrit \emph{kātara}, \textquotedblleft{}confused, discouraged, cowardly\textquotedblright{}; CDIAL lists the Gujarati word, \textquotedblleft{}timid\textquotedblright{} (CDIAL 3027).
\end{cousinweb}

\begin{grammarlens}[title={What you've built}]
One more word for describing what you lived and felt.
\end{grammarlens}

\subsection*{Your turn}
\begin{itemize}
\raggedright
  \item \emph{Say it:} \emph{kāyar}
  \item \emph{Say it:} \emph{kāyar}, once more
  \item \emph{Say it:} \emph{kāyar māṇas jokham joīne bhāgī jāy che.} — A coward runs away at the sight of danger.
  \item \emph{Recall:} say the Gujarati for egotism, then the Gujarati for to get used to, then say \emph{kāyar} again
\end{itemize}

\begin{tcolorbox}[breakable,colback=teal!4,colframe=teal!35!black,title={Before you move on}]
What does \gu{કાયર} mean? (\textquotedblleft{}Cowardly\textquotedblright{}.) Say it once more.
\end{tcolorbox}
\section[\texorpdfstring{\gu{નીડર} (nīḍar)}{નીડર (nīḍar)}]{\gu{નીડર} (nīḍar) — fearless}

\label{lesson:GU-C466-nidar}



\begin{quote}
Before the new one: say the Gujarati for to get used to, then the Gujarati for cowardly.
\end{quote}

\subsection*{You'll want to know: \gu{નીડર}}
\textbf{\gu{નીડર}} — \emph{nīḍar} — \textquotedblleft{}fearless\textquotedblright{}.

\begin{quote}
\gu{નીડર} \gu{મીરા} \gu{અંધારામાં} \gu{એકલી} \gu{ગઈ}.
\end{quote}

\begin{quote}
\emph{nīḍar Mīrā andhārāmā̃ eklī gaī.}
\end{quote}

\begin{quote}
Fearless Mira went out alone in the dark.
\end{quote}

\begin{cousinweb}
\textbf{\gu{નીડર}} (\emph{nīḍar}) is inherited from Sanskrit \emph{nirdara}, \textquotedblleft{}fearless\textquotedblright{}, built from \emph{nir-}, \textquotedblleft{}without\textquotedblright{}, and \emph{dara}, \textquotedblleft{}fear\textquotedblright{}; CDIAL lists the Gujarati word (CDIAL 7339).
\end{cousinweb}

\begin{grammarlens}[title={What you've built}]
One more word for describing what you lived and felt.
\end{grammarlens}

\subsection*{Your turn}
\begin{itemize}
\raggedright
  \item \emph{Say it:} \emph{nīḍar}
  \item \emph{Say it:} \emph{nīḍar}, once more
  \item \emph{Say it:} \emph{nīḍar Mīrā andhārāmā̃ eklī gaī.} — Fearless Mira went out alone in the dark.
  \item \emph{Recall:} say the Gujarati for to get used to, then the Gujarati for cowardly, then say \emph{nīḍar} again
\end{itemize}

\begin{tcolorbox}[breakable,colback=teal!4,colframe=teal!35!black,title={Before you move on}]
What does \gu{નીડર} mean? (\textquotedblleft{}Fearless\textquotedblright{}.) Say it once more.
\end{tcolorbox}
\section[\texorpdfstring{\gu{શૌર્ય} (śaurya)}{શૌર્ય (śaurya)}]{\gu{શૌર્ય} (śaurya) — valour}

\label{lesson:GU-C466-saurya}



\begin{quote}
Before the new one: say the Gujarati for cowardly, then the Gujarati for fearless.
\end{quote}

\subsection*{You'll want to know: \gu{શૌર્ય}}
\textbf{\gu{શૌર્ય}} — \emph{śaurya} — \textquotedblleft{}valour\textquotedblright{}.

\begin{quote}
\gu{આગ} \gu{વખતે} \gu{દાદાએ} \gu{ઘણું} \gu{શૌર્ય} \gu{બતાવ્યું}.
\end{quote}

\begin{quote}
\emph{āg vakhte dādāe ghaṇũ śaurya batāvyũ.}
\end{quote}

\begin{quote}
During the fire my grandfather showed great valour.
\end{quote}

A neuter noun.

\begin{cousinweb}
\textbf{\gu{શૌર્ય}} (\emph{śaurya}) is a tatsama, Sanskrit \emph{śaurya}, \textquotedblleft{}heroism, valour, prowess\textquotedblright{}, as Monier-Williams records it.
\end{cousinweb}

\begin{grammarlens}[title={What you've built}]
One more word for describing what you lived and felt.
\end{grammarlens}

\subsection*{Your turn}
\begin{itemize}
\raggedright
  \item \emph{Say it:} \emph{śaurya}
  \item \emph{Say it:} \emph{śaurya}, once more
  \item \emph{Say it:} \emph{āg vakhte dādāe ghaṇũ śaurya batāvyũ.} — During the fire my grandfather showed great valour.
  \item \emph{Recall:} say the Gujarati for cowardly, then the Gujarati for fearless, then say \emph{śaurya} again
\end{itemize}

\begin{tcolorbox}[breakable,colback=teal!4,colframe=teal!35!black,title={Before you move on}]
What does \gu{શૌર્ય} mean? (\textquotedblleft{}Valour\textquotedblright{}.) Say it once more.
\end{tcolorbox}
\section[\texorpdfstring{\gu{ખમીર} (khamīr)}{ખમીર (khamīr)}]{\gu{ખમીર} (khamīr) — grit; mettle}

\label{lesson:GU-C466-khamir}



\begin{quote}
Before the new one: say the Gujarati for fearless, then the Gujarati for valour.
\end{quote}

\subsection*{You'll want to know: \gu{ખમીર}}
\textbf{\gu{ખમીર}} — \emph{khamīr} — \textquotedblleft{}grit; mettle\textquotedblright{}.

\begin{quote}
\gu{મારા} \gu{દાદામાં} \gu{ઘણું} \gu{ખમીર} \gu{હતું}.
\end{quote}

\begin{quote}
\emph{mārā dādāmā̃ ghaṇũ khamīr hatũ.}
\end{quote}

\begin{quote}
My grandfather had a lot of grit.
\end{quote}

A neuter noun.

\begin{cousinweb}
\textbf{\gu{ખમીર}} (\emph{khamīr}) is usually traced to Arabic \emph{xamīr}, \textquotedblleft{}leaven\textquotedblright{}, through Persian. Gujarati also uses it for the spirit that makes a person rise to a challenge, probably a picture of dough rising.
\end{cousinweb}

\begin{grammarlens}[title={What you've built}]
One more word for describing what you lived and felt.
\end{grammarlens}

\subsection*{Your turn}
\begin{itemize}
\raggedright
  \item \emph{Say it:} \emph{khamīr}
  \item \emph{Say it:} \emph{khamīr}, once more
  \item \emph{Say it:} \emph{mārā dādāmā̃ ghaṇũ khamīr hatũ.} — My grandfather had a lot of grit.
  \item \emph{Recall:} say the Gujarati for fearless, then the Gujarati for valour, then say \emph{khamīr} again
\end{itemize}

\begin{tcolorbox}[breakable,colback=teal!4,colframe=teal!35!black,title={Before you move on}]
What does \gu{ખમીર} mean? (\textquotedblleft{}Grit; mettle\textquotedblright{}.) Say it once more.
\end{tcolorbox}
\section[\texorpdfstring{\gu{ડગમગવું} (ḍagmagvũ)}{ડગમગવું (ḍagmagvũ)}]{\gu{ડગમગવું} (ḍagmagvũ) — to wobble; to waver}

\label{lesson:GU-C466-dagmagvu}



\begin{quote}
Before the new one: say the Gujarati for valour, then the Gujarati for grit.
\end{quote}

\subsection*{You'll want to know: \gu{ડગમગવું}}
\textbf{\gu{ડગમગવું}} — \emph{ḍagmagvũ} — \textquotedblleft{}to wobble; to waver\textquotedblright{}.

\begin{quote}
\gu{તોફાનમાં} \gu{પણ} \gu{તેની} \gu{હિંમત} \gu{ડગમગી} \gu{નહીં}.
\end{quote}

\begin{quote}
\emph{tophānmā̃ paṇ tenī himmat ḍagmagī nahī̃.}
\end{quote}

\begin{quote}
Even in the storm his courage did not waver.
\end{quote}

Past (it …): \gu{તે} \gu{હોડી} \gu{ડગમગી}, \emph{te hoḍī ḍagmagī}.

\begin{cousinweb}
\textbf{\gu{ડગમગવું}} (\emph{ḍagmagvũ}) goes back to a reconstructed \emph{*ḍag-}, \textquotedblleft{}tremble\textquotedblright{}; CDIAL lists the Gujarati verb, \textquotedblleft{}to tremble\textquotedblright{} (CDIAL 5522).
\end{cousinweb}

\begin{grammarlens}[title={What you've built}]
One more word for describing what you lived and felt.
\end{grammarlens}

\subsection*{Your turn}
\begin{itemize}
\raggedright
  \item \emph{Say it:} \emph{ḍagmagvũ}
  \item \emph{Say it:} \emph{ḍagmagvũ}, once more
  \item \emph{Say it:} \emph{tophānmā̃ paṇ tenī himmat ḍagmagī nahī̃.} — Even in the storm his courage did not waver.
  \item \emph{Recall:} say the Gujarati for valour, then the Gujarati for grit, then say \emph{ḍagmagvũ} again
\end{itemize}

\begin{tcolorbox}[breakable,colback=teal!4,colframe=teal!35!black,title={Before you move on}]
What does \gu{ડગમગવું} mean? (\textquotedblleft{}To wobble; to waver\textquotedblright{}.) Say it once more.
\end{tcolorbox}
